% Part: computability
% Chapter: computability-theory
% Section: s-m-n

\documentclass[../../../include/open-logic-section]{subfiles}

\begin{document}

\olfileid[hi]{cmp}{thy}{smn}
\olsection{$s$-$m$-$n$ प्रमेय}

\begin{explain}
अगला प्रमेय ``$s$-$m$-$n$ प्रमेय'' कहलाता है; इसका कारण अभी स्पष्ट हो
जाएगा। कठिन भाग यह समझना है कि प्रमेय कहता क्या है; कथन समझ में आ जाने
पर वह काफ़ी स्पष्ट प्रतीत होगा।
\end{explain}

\begin{thm}
\ollabel{thm:s-m-n}
  प्राकृतिक संख्याओं $n$ और~$m$ के प्रत्येक युग्म के लिए एक आदिम
  पुनरावर्ती फलन~$s^m_n$ है, ऐसा कि प्रत्येक अनुक्रम
  $e$, $a_0$, \dots, $a_{m-1}$, $y_0$ ,\dots, $y_{n-1}$ के लिए हमें
  \[
  \cfind{s^m_n(e, a_0, \dots, a_{m-1})}[n](y_0, \dots, y_{n-1}) \simeq
  \cfind{e}[m+n](a_0, \dots, a_{m-1}, y_0, \dots, y_{n-1}).
\]
\end{thm}

\begin{explain}
$s^m_n$ को \emph{प्रोग्रामों} पर क्रिया करने वाला मानना उपयोगी है।
अर्थात् $s^m_n$, प्रोग्राम~$e$ लेता है जो किसी $(m+n)$-स्थानी फलन के लिए
है, और साथ में नियत आगत $a_0$, \dots, $a_{m-1}$ लेता है; तथा प्रोग्राम
$s^m_n(x, a_0, \dots, a_{m-1})$ लौटाता है, जो शेष तर्कों के $n$-स्थानी
फलन के लिए है।
\iftag{TMs}{यदि $x$ को किसी ट्यूरिंग मशीन का वर्णन माना जाए, तो
$s^m_n(e, a_0, \dots, a_{m-1})$ वह ट्यूरिंग मशीन है जो आगत $y_0$,
\dots,~$y_{n-1}$ मिलने पर आगत शृंखला के आरंभ में $a_0$, \dots,
~$a_{m-1}$ लगाती है और~$e$ चलाती है। तब प्रत्येक $s^m_n$ केवल वह आदिम
पुनरावर्ती फलन है जो उपयुक्त ट्यूरिंग मशीन का कूट ज्ञात करता है।}{}
\end{explain}

\end{document}

